\chapter{共分散の推定}\label{ch:cov}

平均分散最適化も Black--Litterman モデルも、共分散の行列 $\Sigma$ を入力にとる。$\Sigma$ は観測できないので、過去のリターンから推定する。この章では、推定の難しさと、その対処の方法を説明する。補足実験 X01〜X04 は、この章の方法を比べたものである。

\section{標本共分散と、その問題}

過去 $T$ 日の $N$ 銘柄の日次リターンを並べた $T\times N$ の行列を $X$（各列から平均を引いたもの）とすると、\term{標本共分散}は
\begin{equation}
S = \frac{1}{T-1}X\T X
\end{equation}
である。この推定には 2 つの問題がある。

\begin{description}[leftmargin=1.5em,style=nextline,itemsep=0.4em]
\item[推定するものが多すぎる] $N$ 銘柄の共分散の行列には $N(N+1)/2$ 個の値がある。$N=500$ なら 125,250 個である。1 年（$T=252$ 日）のデータには $252\times500=126{,}000$ 個の数しかない。値の数とデータの数がほぼ同じでは、1 つ 1 つの推定はとても粗い。
\item[特異になる] $S$ の階数は高々 $T-1$ である。$N > T$ なら、$S$ は逆行列を持たない（固有値の多くが 0 になる）。平均分散最適化は $\Sigma^{-1}$ を使うので、そのままでは計算できない。$N<T$ でも $N/T$ が 1 に近いと、最小の固有値がとても小さくなり、逆行列がひどく不安定になる。
\end{description}

\subsection{最適化が誤差を増幅する}

平均分散最適化は、リスクが小さく見える組み合わせに大きく賭ける。推定の誤差で「たまたま相関が低く見えた」銘柄の組や、「たまたまリスクが小さく見えた」方向に、比率が集まる。Michaud はこれを\term{誤差の最大化}と呼んだ。最適化の結果が、推定の誤差に最も敏感な方向に偏るのである。

\subsection{固有値で見る}

共分散の行列を固有値に分解すると $S = \sum_k \lambda_k v_k v_k\T$ になる。大きい固有値の方向は市場全体やセクターの動きで、小さい固有値の方向は「ほとんど動かない組み合わせ」である。最適化は小さい固有値の方向を好む（そこが低リスクに見えるため）。ところが、小さい固有値ほど推定の誤差が相対的に大きい。\term{条件数} $\lambda_{\max}/\lambda_{\min}$ は、この不安定さの目安である。

\section{ランダム行列の理論：Marcenko--Pastur の分布}

銘柄の間に本当は何の相関もなく、リターンが分散 $\sigma^2$ の独立な雑音だけだったとしても、標本共分散の固有値はばらつく。$N,T\to\infty$、$q=N/T$ を一定にしたとき、固有値の分布は\term{Marcenko--Pastur の分布}に近づき、その範囲は
\begin{equation}
\lambda_{\pm} = \sigma^2\left(1 \pm \sqrt{q}\right)^2
\end{equation}
になる。$q=457/252\approx1.8$ なら、$\lambda_+ \approx 5.5\sigma^2$ である（$q>1$ のときは、これとは別に $N-T$ 個以上の固有値が 0 になる）。標本の固有値のうち $\lambda_+$ より小さいものは、雑音と区別できない。

\begin{defbox}[定義：固有値のノイズ除去（López de Prado）]
相関の行列を固有値に分解し、$\lambda_+$ より小さい固有値をすべて、その平均で置き換えてから行列を組み立て直す。大きい固有値（市場やセクターの動き）は残し、雑音の部分だけを平らにする。対角を 1 に戻し、各銘柄の分散を掛けて共分散に戻す。
\end{defbox}

\section{縮小推定：Ledoit--Wolf}

\term{縮小推定}は、標本共分散 $S$ を、単純で安定した\term{目標} $F$ に向けて縮める。
\begin{equation}
\hat\Sigma = \delta F + (1-\delta) S, \qquad 0\le\delta\le1
\end{equation}
目標には、単位行列の定数倍（すべての銘柄が同じ分散で無相関）や、相関がすべて等しい行列を使う。$S$ は偏りがないが分散が大きく、$F$ は偏っているが分散が小さい。両者を混ぜると、平均二乗誤差が最も小さくなる $\delta$ がある。Ledoit と Wolf は、その $\delta$ をデータから推定する式を導いた。$\hat\Sigma$ は常に正定値で、逆行列が安定する。

縮小推定は、固有値で見ると「大きい固有値を小さく、小さい固有値を大きく」して、固有値の広がりを縮めている。Marcenko--Pastur のノイズ除去と目的は同じで、やり方が違う。

\section{窓の長さと指数加重}

過去のどこまでを使うか（\term{窓}）も選ぶ必要がある。

\begin{itemize}
\item \textbf{短い窓}（63 日など）：最近のリスクの変化に早く追従するが、推定の誤差が大きい。
\item \textbf{長い窓}（504 日など）：推定は安定するが、市場の状態が変わっても古いリスクを引きずる。
\end{itemize}

\term{指数加重移動平均}（EWMA）は、最近のデータほど重みを大きくする。$k$ 日前のデータの重みを $\lambda^k$（$0<\lambda<1$）にし、\term{半減期} $h$ を使って $\lambda = 0.5^{1/h}$ とする。半減期 63 日なら、63 日前のデータの重みは今日の半分、126 日前は 4 分の 1 になる。RiskMetrics の日次の標準は $\lambda=0.94$（半減期約 11 日）である。

\section{リスクの予測を評価する}

共分散は\textbf{将来のリスクの予測}として使うので、その良し悪しは、予測したリスクと実際に実現したリスクを比べて評価する。

\begin{itemize}
\item \textbf{比}：ポートフォリオの予測の標準偏差 $\hat\sigma_p = \sqrt{w\T\hat\Sigma w}$ と、実現した標準偏差 $\sigma_p^{\text{real}}$ の比。1 に近いほどよい。最適化したポートフォリオでは、比は 1 より大きくなりがちである（最適化が「リスクが小さく見える方向」を選ぶので、そこで予測が過小になる）。
\item \textbf{QLIKE}：分散の予測 $h$ と実現 $r^2$ について $\text{QLIKE} = r^2/h - \log(r^2/h) - 1$ を平均した損失。分散の予測の比較によく使われる。
\end{itemize}

\begin{notebox}[実現の TE は約束より大きくなる]
最適化したポートフォリオの予測のリスクは、実現より小さく出る。このプロジェクトでは、TE（\chapref{ch:active}）を 3\% に抑えるよう最適化したのに、実現の TE は平均 4.6〜4.8\% だった（X02）。窓や EWMA を変えても 1.4 倍のままで、過去 36 か月の「実現 ÷ 予測」の比で上限を割る校正（X03）で、やっと 3.4\% 前後になった。
\end{notebox}

\section{そのほかの方法}

\begin{description}[leftmargin=1.5em,style=nextline]
\item[ファクター・モデル] $\Sigma = B\,\Sigma_f\,B\T + D$。少数のファクター（市場、セクター、規模など）の共分散 $\Sigma_f$ と、各銘柄のファクターへの感応度 $B$、個別の分散 $D$（対角）で表す。推定するものが $N^2$ から $N\times K$ 程度に減る。商用のリスク・モデル（Barra、Axioma）はこの形である。
\item[GARCH] 分散が時間とともに変わることを、過去の分散と過去の衝撃の式で表す。1 銘柄には強いが、500 銘柄の行列に広げるのは難しい。
\item[予想ボラティリティで合わせる] 行列の形は過去から推定し、全体の大きさだけを VIX などの市場の予想に合わせる（X04）。
\end{description}

\begin{projbox}
\begin{itemize}
\item \textbf{今の BLM}：過去 252 取引日の日次リターンの Ledoit--Wolf の縮小推定。
\item \textbf{X01}：標本共分散、Ledoit--Wolf、Marcenko--Pastur のノイズ除去を、2016〜2026 年の 127 か月で比べた。ノイズ除去が勝つのは制約のない最小分散だけで、BLM の最適化（TE の上限つき）では Ledoit--Wolf で十分だった。
\item \textbf{X02}：窓 63〜504 日と EWMA（半減期 63・126 日）。ewma126 が最良だったが、実現の TE は約束の約 1.4 倍のまま。
\item \textbf{X03}：上限の校正。実現の TE 4.79\% → 3.38\%。実験 7・9 で使った。
\item \textbf{X04}：VIX で大きさを合わせても、校正に足す改善はなかった。
\end{itemize}
\end{projbox}
